% Q4: Ideal transformer impedance matching — Vs + Rs → transformer → RL
\documentclass[border=10pt]{standalone}
\usepackage[american voltages, american currents]{circuitikz}
\usepackage{amsmath}
\renewcommand{\familydefault}{\sfdefault}

\begin{document}
\begin{circuitikz}[
    line width=0.8pt,
    every node/.style={font=\sffamily}
]

% Voltage source (left)
\draw (0,4) to[V, v=$V_s$] (0,0);

% Series resistance Rs
\draw (0,4) to[R, l=$R_s$] (3,4);

% Primary terminal dots
\fill (3,4) circle (2pt);
\fill (3,0) circle (2pt);

% Ideal transformer
% Primary winding
\draw (3,4) to[L] (3,1.5);
\draw (3,1.5) -- (3,0);

% Core lines
\draw[thick] (3.8,0.5) -- (3.8,3.7);
\draw[thick] (4.2,0.5) -- (4.2,3.7);

% Secondary winding
\draw (5,4) to[L] (5,1.5);
\draw (5,1.5) -- (5,0);

% Turns ratio label
\node[above] at (4,4.3) {$N_1 : N_2$};

% Dots on transformer
\fill (3.2,3.8) circle (3pt);
\fill (4.8,3.8) circle (3pt);

% Ideal label
\node at (4,2) {\footnotesize ideal};

% Secondary terminal dots
\fill (5,4) circle (2pt);
\fill (5,0) circle (2pt);

% Load resistor
\draw (5,4) -- (7,4);
\draw (7,4) to[R, l=$R_L$] (7,0);
\draw (7,0) -- (5,0);

% Bottom return wire (source side)
\draw (0,0) -- (3,0);

% Zin arrow at primary terminals
\draw[<->, >=stealth, thick, red!70!black] (2.6,3.5) -- (2.6,0.5);
\node[left, red!70!black] at (2.5,2) {$Z_{\text{in}}$};

% Current arrow
\draw[->, >=stealth, thick] (1.0,4.4) -- (2.5,4.4)
    node[midway, above] {$I_1$};

\end{circuitikz}
\end{document}
